%==========================================================================================
\section{\texttt{xrf\_fit}: the spectrum model and the fits}\label{sec:fit}
%==========================================================================================

\subsection{Overview: the constrained model}

An energy-dispersive XRF spectrum is a sum of peaks on a continuum. Fitting every peak with
free position, width and height is unstable as soon as peaks overlap, and produces numbers
that cannot be compared between elements. \fn{xrf_fit} instead builds a \emph{constrained}
model in which almost everything except one amplitude per element and shell is computed:
\begin{equation}
  y(c) = B(c) + \sum_g A_g \sum_{l\in g} r_l\,\tau_l(d)
  \Big[\big(1-\textstyle\sum_e f_{le}\big)\,R\big(c-c_l;\sigma_l\big)
       + \sum_e f_{le}\,R\big(c-c_{le};\sigma_{le}\big)\Big]
  + Y_\mathrm{scatter}(c) + \tau_\mathrm{pp}\,P(c),
  \label{eq:model}
\end{equation}
where
\begin{itemize}[nosep]
\item $c$ is the MCA channel, and $E(c) = Gc+O$ the energy calibration;
\item $B(c)$ is a polynomial background (or the exponential of a polynomial);
\item $g$ runs over \emph{line groups}, one per (element, shell family): \fn{"Cu K"},
  \fn{"Pb L"}, \fn{"Pb M"};
\item $A_g$, the only free amplitude of the group, is the integrated number of counts of its
  most intense line (the \emph{reference line});
\item $r_l$ is the intensity of line $l$ relative to the reference line ($r=1$ for it),
  computed for the excitation, the detector and, optionally, the sample matrix;
\item $\tau_l(d)$ rescales $r_l$ when the sample--detector air path $d$ is fitted;
\item $c_l = (E_l-O)/G$ is the position of the line, from its tabulated energy; $\sigma_l$ its
  Gaussian width from the detector resolution law;
\item $f_{le}$ and $c_{le}$ are the fraction and position of its sensor escape peaks;
\item $R$ is the unit-area detector response (Gaussian or Voigt, optionally with a
  low-energy tail and a step);
\item $Y_\mathrm{scatter}$ holds the elastic and Compton peaks under monochromatic
  excitation, and $\tau_\mathrm{pp}P(c)$ the pile-up spectrum.
\end{itemize}
An $N$-element fit has $N$ amplitudes plus a handful of shared parameters, instead of several
per peak. Peak positions are never free individually: they move only through the
calibration. This is the approach of the fundamental-parameter fitting codes such as PyMca
\citep{Sole2007}, written here for white-beam excitation and for maps.

The module also provides the statistical treatment of map data (normalisation, averaging,
uncertainties), the per-pixel linear solve that turns the average-spectrum fit into maps with
error maps, the tracking of detector-resolution changes during a scan, and diagnostic plots.

\pycode{xrf_fit}{_header}

\subsection{Energy calibration}

\api[\texttt{Calibration}]{\texttt{Calibration(gain=10.032, offset=-31.263)}}
Linear calibration $E\,[\mathrm{eV}] = G\,c + O$ (a frozen dataclass). The defaults were
measured on Ga/As K lines with the BM32 XIA configuration of 2026; \fn{xrf_calibration}
measures it on reference samples.
\begin{args}
\item[gain] $G$, eV per channel.
\item[offset] $O$, energy of channel 0 in eV.
\end{args}
\blk{Methods} \fn{energy(channel)} and \fn{channel(energy_ev)} (not rounded) convert in
either direction, element-wise; \fn{channel_range(emin_ev, emax_ev, n_channels)} returns the
half-open channel range $[\lfloor c(E_\mathrm{min})\rfloor, \lceil c(E_\mathrm{max})\rceil+1)$
clipped to the spectrum.
\pycode{xrf_fit}{Calibration}
\pycode{xrf_fit}{Calibration.energy}
\pycode{xrf_fit}{Calibration.channel}
\pycode{xrf_fit}{Calibration.channel_range}

\subsection{The detector}\label{sec:detector}

\api[\texttt{Detector}]{\texttt{Detector(material="Si", noise\_ev=30.0, fano=0.12, gamma\_ev=5.0, sensor\_thickness\_um=450.0, window\_material="C", window\_density\_g\_cm3=2.2, window\_thickness\_um=1.0, air\_path\_cm=20.0, air\_temperature\_c=20.0, air\_pressure\_hpa=1013.25, dead\_layer\_um=0.0, tail\_fraction=0.0, tail\_slope=1.5, step\_fraction=0.0)}}
The energy resolution, detection efficiency and line shape of a solid-state detector. The
defaults describe the BM32 KETEK VITUS H7 silicon drift detector \citep{Lechner1996}: 450\,\textmu m
of Si behind a 1\,\textmu m carbon window, about 20\,cm from the sample.
\begin{args}
\item[material] sensor, \fn{"Si"} or \fn{"Ge"}: sets the pair-creation energy $\varepsilon$
  (3.85 and 2.96\,eV) and the density.
\item[noise\_ev] electronic-noise contribution $N$ to $\sigma$.
\item[fano] Fano factor $F$.
\item[gamma\_ev] Lorentzian half width $\Gamma$ of the Voigt profile.
\item[sensor\_thickness\_um] active thickness $d_s$.
\item[window\_*] entrance-window material, density and thickness.
\item[air\_path\_cm, air\_temperature\_c, air\_pressure\_hpa] sample--window air path and
  ambient conditions.
\item[dead\_layer\_um] inactive Si at the entrance.
\item[tail\_fraction, tail\_slope, step\_fraction] line-shape parameters, see \fn{peak_shape}.
\end{args}
\blk{Resolution} The Gaussian width combines electronic noise and the statistics of charge
creation \citep{Fano1947,Knoll2010}:
\begin{equation}
  \sigma^2(E) = N^2 + F\varepsilon E, \qquad \mathrm{FWHM} = 2\sqrt{2\ln2}\,\sigma =
  2.3548\,\sigma .
  \label{eq:resolution}
\end{equation}
The defaults ($N = 30$\,eV, $F = 0.12$) give 142\,eV FWHM at 5.9\,keV (the class docstring
quotes ``$\sim$135\,eV'', the middle of the 125--150\,eV range of the BM32 detector; the
resolution is fitted in practice).
\begin{hyp}
Only the product $F\varepsilon$ enters Eq.~\eqref{eq:resolution}. The value of $\varepsilon$
for Si used here, 3.85\,eV, is above the measured 3.66\,eV at room temperature
\citep{Scholze1998}; since $F$ is fitted, this only rescales the fitted Fano factor and not
the widths. The Lorentzian width $\Gamma$ is physically the natural width of the level (a few
eV); a fitted value above $\sim$30\,eV signals that a tail term is needed.
\end{hyp}
\pycode{xrf_fit}{Detector}

\api[\texttt{pair\_energy\_ev}]{\texttt{pair\_energy\_ev} (property)}
$\varepsilon$ for the sensor material; \fn{ValueError} for other materials.
\pycode{xrf_fit}{Detector.pair_energy_ev}

\api[\texttt{Detector.efficiency}]{\texttt{Detector.efficiency(energy\_ev)}}
Probability that a photon emitted towards the detector is recorded in the full-energy
region:
\begin{equation}
  \epsilon(E) = \underbrace{\e^{-\mu_\mathrm{air}(E)\rho_\mathrm{air}d_\mathrm{air}}}_{\text{air}}\;
  \underbrace{\e^{-\mu_w(E)\rho_w t_w}}_{\text{window}}\;
  \underbrace{\e^{-\mu_\mathrm{Si}(E)\rho_\mathrm{Si}t_\mathrm{dead}}}_{\text{dead layer}}\;
  \underbrace{\left[1-\e^{-\tau_\mathrm{Si}(E)\rho_\mathrm{Si}d_s}\right]}_{\text{absorbed}} .
  \label{eq:efficiency}
\end{equation}
Attenuation along the path uses total cross-sections (any interaction removes the photon);
absorption in the sensor uses the photoelectric cross-section $\tau$ only, since a scattered
photon does not deposit its full energy. Escape peaks are not subtracted: they redistribute
counts within the recorded spectrum and belong to the line shape. Returns a float for a
scalar input, else an array; energies where xraylib has no data give 0.
\begin{hyp}
Normal incidence on a uniform sensor, collimator and grid ignored, single air-path length.
The air path dominates the efficiency of soft lines: 20\,cm of air transmit about 4\,\% at
In~L$\alpha$ (3.29\,keV), and each extra centimetre costs about 15\,\% there, against 0.3\,\%
at Ga~K$\alpha$.
\end{hyp}
\pycode{xrf_fit}{Detector.efficiency}

\api[\texttt{Detector.escape\_lines}]{\texttt{Detector.escape\_lines(energy\_ev)}}
Sensor escape peaks of a line. A photon absorbed in the K shell of the sensor can be followed
by a sensor K X-ray that leaves the crystal, so the event is recorded at $E-E_X$. For normal
incidence on a semi-infinite sensor \citep{ReedWare1972},
\begin{equation}
  f_X = \tfrac12\,\omega_K\,\frac{\tau_K}{\tau}\,p_X
  \left[1-\frac{\mu_X}{\mu_E}\ln\!\left(1+\frac{\mu_E}{\mu_X}\right)\right],
  \label{eq:escape}
\end{equation}
with $\omega_K$ the K fluorescence yield of the sensor, $\tau_K/\tau$ the K-shell share of
photoabsorption at $E$, $p_X$ the radiative branching of the escaping line (K$\alpha$ or
K$\beta$), $\mu_E$ the photoabsorption coefficient at $E$ and $\mu_X$ the attenuation of the
escaping X-ray. Returns \fn{[(escape energy in eV, fraction), ...]}, empty below the sensor's
K edge. Typical values: 0.1\,\% for Ga/As K in Si, about 1\,\% just above the Si K edge.
\pycode{xrf_fit}{Detector.escape_lines}

\api[\texttt{sigma\_ev}, \texttt{fwhm\_ev}, \texttt{from\_fwhm}]{\texttt{sigma\_ev(energy\_ev, noise\_ev=None, fano=None)}, \texttt{fwhm\_ev(energy\_ev, **kw)}, \texttt{from\_fwhm(fwhm\_ev, at\_energy\_ev, material="Si", fano=0.12, gamma\_ev=5.0, **kw)} (class method)}
$\sigma(E)$ and FWHM from Eq.~\eqref{eq:resolution}, optionally with other noise and Fano
values (used by the fit). \fn{from_fwhm} solves Eq.~\eqref{eq:resolution} for $N$ given a
quoted resolution (``135\,eV at Mn K$\alpha$'') and a Fano factor, and raises
\fn{ValueError} if the FWHM is below the Fano limit; extra keywords set the efficiency fields.
\pycode{xrf_fit}{Detector.sigma_ev}
\pycode{xrf_fit}{Detector.fwhm_ev}
\pycode{xrf_fit}{Detector.from_fwhm}

\subsection{Fluorescence lines and line groups}\label{sec:lines}

\api[\texttt{Line}, \texttt{LineGroup}]{\texttt{Line(name, energy\_ev, rel)}, \texttt{LineGroup(element, family, lines, ref\_strength=nan)}}
A \fn{Line} is an IUPAC name without the \fn{_LINE} suffix, a tabulated energy and an
intensity relative to the reference line of its group. A \fn{LineGroup} is a set of lines of
one element and one shell family (\fn{"K"}, \fn{"L"}, \fn{"M"}) sharing one amplitude;
\fn{lines[0]} is the reference line. \fn{ref_strength} keeps the absolute strength of the
reference line---the normalisation that \fn{rel} divides out: cross-section (cm$^2$/g) at
the excitation energy, or flux-weighted cross-section, or matrix sensitivity, times the
detection efficiency if a detector was given. It converts fitted areas into atom ratios.
Properties: \fn{label} (\fn{"Cu K"}), \fn{reference}, and \fn{total_rel} ($\sum_l r_l$,
converting the reference-line area into the area of the whole modelled group).
\pycode{xrf_fit}{Line}
\pycode{xrf_fit}{LineGroup}
\pycode{xrf_fit}{LineGroup.label}
\pycode{xrf_fit}{LineGroup.reference}
\pycode{xrf_fit}{LineGroup.total_rel}

\api[\texttt{\_family\_of}, \texttt{\_all\_line\_macros}, \texttt{\_is\_iupac\_line}]{\texttt{\_family\_of(macro\_name)}, \texttt{\_all\_line\_macros()}, \texttt{\_is\_iupac\_line(macro\_name)}}
xraylib exposes every transition twice---under its Siegbahn name (\fn{KA1}, \fn{LB2},
\fn{LG1}, \fn{LL}) and under its IUPAC name (\fn{KL3}, \fn{L3N5}, \fn{L2N4}, \fn{L3M1})---plus
composite macros (\fn{KA} = KL3 + KL2). Enumerating all of them would count intensity two or
three times, so only the IUPAC set is used: it is complete, non-redundant, and has one name
per transition. The regular expression demands a full subshell designator on both sides
(\fn{L3M5} matches, \fn{LL} does not). \fn{_family_of} gives the family from the initial
vacancy; \fn{_all_line_macros} collects every \fn{*_LINE} constant of xraylib.
\pycode{xrf_fit}{_family_of}
\pycode{xrf_fit}{_all_line_macros}
\pycode{xrf_fit}{_is_iupac_line}

\api[\texttt{\_line\_energy\_ev}]{\texttt{\_line\_energy\_ev(z, macro)}}
Tabulated line energy in eV \citep{Elam2002,Schoonjans2011}, or \fn{None} if the line does
not exist for that element.
\pycode{xrf_fit}{_line_energy_ev}

\api[\texttt{\_line\_strength}]{\texttt{\_line\_strength(z, macro, e0\_kev, weighting)}}
Unnormalised intensity of a line under monochromatic excitation at $E_0$:
\begin{itemize}[nosep]
\item \fn{"kissel"}: the partial fluorescence cross-section $Q_l(E_0)$ (cm$^2$/g), which
  includes the photo-ionisation cross-section of each subshell, the Coster--Kronig transfers
  between subshells and the fluorescence yields \citep{Krause1979,Campbell2003}, and is
  therefore comparable across subshells;
\item \fn{"radrate"}: the radiative rate, normalised \emph{within} each subshell, hence
  comparable only between lines sharing an initial vacancy. Across an L series it is wrong:
  for Ir it predicts L$\beta_1$/L$\alpha_1$ = 1.13 where the Kissel cross-sections give 0.65
  and the measured ratio on the NIO295 map is 0.77.
\end{itemize}
Returns 0 below the edge.
\pycode{xrf_fit}{_line_strength}

\api[\texttt{\_line\_strength\_spectral}]{\texttt{\_line\_strength\_spectral(z, macro, energy\_kev, flux)}}
The correct strength under polychromatic excitation:
\begin{equation}
  S_l = \int Q_l(E)\,\frac{\dd\Phi}{\dd E}\,\dd E ,
  \label{eq:flux-weighted}
\end{equation}
with the per-0.1\,\% bandwidth flux converted to a density per keV
(Eq.~\eqref{eq:per-energy}; omitting the Jacobian inflated the Ga/As k-factor by about 9\,\%
on BM32). Each subshell is excited only by the part of the spectrum above its own edge. The
integral uses the trapezoidal rule, which cannot overshoot at the edge discontinuity
(Simpson's rule can); with a grid step $h$ the edge error is of order $h$ times the jump, so a
grid of $\sim$10\,eV keeps it below 0.1\,\%.
\pycode{xrf_fit}{_line_strength_spectral}

\api[\texttt{element\_line\_groups}]{\texttt{element\_line\_groups(elements, emin\_ev, emax\_ev, e0\_eff\_kev=20.0, families=("K","L","M"), min\_rel=0.01, weighting="kissel", merge\_within\_ev=0.0, incident\_spectrum=None, detector=None, matrix=None)}}
Builds the line groups to fit.
\begin{args}
\item[elements] symbols or atomic numbers.
\item[emin\_ev, emax\_ev] only lines inside this window are kept (normally the fit range).
\item[e0\_eff\_kev] single excitation energy for the cross-sections; it must exceed the
  highest relevant edge, or high-subshell lines silently get zero intensity.
\item[families] shell families considered.
\item[min\_rel] lines weaker than this fraction of the group's reference line are dropped.
\item[weighting] \fn{"kissel"} or \fn{"radrate"} (ignored with a spectrum or a matrix).
\item[merge\_within\_ev] if positive, lines closer than this are merged into one effective
  line at their intensity-weighted energy (to shrink the model; the detector cannot resolve
  them anyway).
\item[incident\_spectrum] \texttt{(energy\_kev, flux\_per\_0.1\%BW)}: strengths become
  Eq.~\eqref{eq:flux-weighted}.
\item[detector] multiplies every strength by $\epsilon(E_l)$, so that relative intensities
  and \fn{ref_strength} describe the \emph{recorded} spectrum.
\item[matrix] an \fn{xrf_matrix.MatrixModel}: strengths include self-absorption (and
  secondary fluorescence for the lines that can survive \fn{min_rel}); it carries its own
  spectrum, so \fn{incident_spectrum} and \fn{e0_eff_kev} are ignored. \fn{ref_strength}
  then means sensitivity per unit mass fraction in the target layer.
\end{args}
\blk{Returns} one \fn{LineGroup} per (element, family) with at least one line left, in the
order of \fn{elements}; lines sorted by decreasing relative intensity.
\blk{Method} For each element and family, every IUPAC line of that family inside the window
gets its strength (cross-section, flux-weighted cross-section, or matrix emission), times
the efficiency; the strongest defines the reference; lines below \fn{min_rel} are dropped.
\pycode{xrf_fit}{element_line_groups}

\api[\texttt{\_merge\_close\_lines}, \texttt{\_collapse}]{\texttt{\_merge\_close\_lines(found, tol\_ev)}, \texttt{\_collapse(bucket)}}
Lines sorted by energy are grouped while consecutive lines are within \fn{tol_ev}; each group
becomes one line at the strength-weighted mean energy, with the summed strength and the name
of its strongest member plus the number of merged lines (\fn{"L3M5+1"}).
\pycode{xrf_fit}{_merge_close_lines}
\pycode{xrf_fit}{_collapse}

\subsection{Normalisation, averaging and counting statistics}

\api[\texttt{normalise\_spectra}]{\texttt{normalise\_spectra(spectra, monitor=None, live\_time=None, normalise\_monitor=True)}}
Corrects raw MCA spectra for the incident flux and the detector dead time,
\begin{equation}
  y_i = \frac{x_i}{s_i}, \qquad s_i = \frac{M_i}{\langle M\rangle}\,t_{\mathrm{live},i},
\end{equation}
with $M_i$ the monitor of pixel $i$ (divided by its mean if \fn{normalise_monitor}) and
$t_\mathrm{live}$ the live time reported by the XIA electronics (\fn{xia_det0_live_time}).
Dividing by the live time turns recorded counts into an incoming rate. On BM32 white-beam
maps the dead-time fraction varies strongly with position (0.33--0.79 on the NIO295 map, a
factor 1.5--4.8 in the correction) because it follows the total scattering power, so it does
not cancel in elemental ratios.
\blk{Returns} \fn{(corrected, scale)}; pixels with a non-positive or non-finite scale are NaN.
\pycode{xrf_fit}{normalise_spectra}

\api[\texttt{\_pixel\_rows}]{\texttt{\_pixel\_rows(spectra, name)}}
Reshapes any array whose last axis is the channel axis---a single spectrum, a list of
pixels, a gridded map---into $(n_\mathrm{pixels}, n_\mathrm{channels})$.
\pycode{xrf_fit}{_pixel_rows}

\api[\texttt{average\_spectrum}]{\texttt{average\_spectrum(corrected)}}
The mean of the valid pixels (those without NaN) and the standard error of that mean,
$s/\sqrt n$, per channel. The standard error contains both counting noise and real
pixel-to-pixel variation. With a single spectrum (\fn{ct} scan) the standard error is
undefined and returned as NaN; use \fn{counting_error} for the fit weights.
\blk{Returns} \fn{(mean, sem, n_used)}.
\pycode{xrf_fit}{average_spectrum}

\api[\texttt{counting\_error}]{\texttt{counting\_error(raw\_spectra, scale)}}
The Poisson uncertainty of the normalised average. With $\bar y = \frac1n\sum_i x_i/s_i$ and
$\mathrm{Var}(x_i) = x_i$,
\begin{equation}
  \sigma_{\bar y} = \frac1n\sqrt{\sum_i \frac{x_i}{s_i^2}} .
\end{equation}
This is the right weight for fitting the \emph{average} spectrum: pixel-to-pixel variation is
signal, not noise, as far as the average is concerned, and on a heterogeneous map the
standard error can exceed this by an order of magnitude. Accepts any scan shape; \fn{scale}
can be a scalar.
\pycode{xrf_fit}{counting_error}

\subsection{Pile-up}

\api[\texttt{pileup\_spectra}]{\texttt{pileup\_spectra(raw\_spectra, calibration, real\_time=1.0, chunk=512)}}
Per-pixel pile-up spectra per second of resolving time. For random arrivals, two events
closer than the pulse-pair resolving time $\tau$ are recorded as one event at $E_1+E_2$. With
$x$ the measured spectrum, recorded in the real time $t$, the number of such pairs at sum
energy $E$ is \citep{Statham1977}
\begin{equation}
  N_\mathrm{pp}(E) = \frac{\tau}{t}\,(x\otimes x)(E),
\end{equation}
where $\otimes$ is the self-convolution over ordered pairs; the total, $\tau N^2/t =
N\,r\,\tau$ with $r$ the count rate, is the expected number of pairs. The function returns
$(x\otimes x)/t$ on the MCA channel grid, so multiplying by $\tau$ in seconds gives counts.
\blk{Method} The convolution is computed with FFTs (padded to a power of two above $2n$,
\citealp{Cooley1965}), in batches of \fn{chunk} pixels. The self-pair term is made unbiased
by subtracting $x_c$ at index $2c$ ($x_c(x_c-1)$ pairs instead of $x_c^2$), which matters at
low counts. Index $n$ of the convolution has energy $Gn+2O$ (the offset counts twice), so it
is mapped back to channel $c = n + O/G$ by linear interpolation.
\begin{hyp}
Random (Poisson) arrivals and a single resolving time; pile-up of three or more events is
neglected. Pile-up removes parent events with a probability independent of their energy, so
it does not bias line \emph{ratios}; it adds sum peaks that other components would otherwise
absorb.
\end{hyp}
\pycode{xrf_fit}{pileup_spectra}

\api[\texttt{average\_pileup\_template}]{\texttt{average\_pileup\_template(pileup, scale)}}
Pile-up template for the normalised average spectrum, per microsecond of resolving time:
$10^{-6}\,\langle P_i/s_i\rangle$. It must be built from the per-pixel spectra, since pile-up
goes as the square of the local rate; the fitted amplitude is then $\tau$ in \textmu s.
\pycode{xrf_fit}{average_pileup_template}

\subsection{Background estimate}

\api[\texttt{snip\_background}]{\texttt{snip\_background(y, width=40, iterations=24, smooth=True)}}
The Statistics-sensitive Non-linear Iterative Peak-clipping algorithm \citep{Ryan1988}. The
spectrum (optionally smoothed over 3 channels) is transformed by the log-log-square-root
operator $v = \ln[\ln(\sqrt{y+1}+1)+1]$, which compresses peaks and makes the clipping
statistics-aware; then each point is replaced by
$v_c \leftarrow \min\!\left(v_c, \tfrac12(v_{c-w}+v_{c+w})\right)$ for windows $w = 1,\dots,
W$ and then decreasing again for \fn{iterations} passes (decreasing windows smooth the result,
\citealp{Morhac1997}); finally the transform is inverted. \fn{width} should exceed the
half-width of the widest peak. Used by the coarse calibration; the fits themselves model
the background.
\pycode{xrf_fit}{snip_background}

\subsection{Line shape}

\api[\texttt{peak\_shape}]{\texttt{peak\_shape(x, sigma, gamma=0.0, tail\_fraction=0.0, tail\_slope=1.5, step\_fraction=0.0)}}
The unit-area detector response to one line, plus its step, of the Hypermet type
\citep{Phillips1976,VanGysel2003}:
\begin{equation}
  R(x) = (1-t)\,V(x;\sigma,\Gamma)
  + \frac{t}{2\beta}\,\exp\!\left(\frac{x}{\beta}+\frac{\sigma^2}{2\beta^2}\right)
    \erfc\!\left(\frac{x}{\sqrt2\sigma}+\frac{\sigma}{\sqrt2\beta}\right)
  + \frac{s}{2\sqrt{2\pi}\sigma}\,\erfc\!\left(\frac{x}{\sqrt2\sigma}\right),
  \label{eq:hypermet}
\end{equation}
where $x$ is the distance to the line centre (in channels in the model), $V$ is the Voigt
profile (\fn{scipy.special.voigt_profile}, $\Gamma$ the Lorentzian half width) or the Gaussian
when $\Gamma = 0$, $t$ the tail fraction, $\beta = $ \fn{tail_slope}$\times\sigma$ the decay
length of the tail, and $s$ the step height relative to the Gaussian peak height
$1/(\sqrt{2\pi}\sigma)$. The tail term is the convolution of a one-sided exponential
$\e^{x/\beta}/\beta$ ($x<0$) with the Gaussian, and has unit area; the step is the
convolution of a flat shelf with the Gaussian. Both describe incomplete charge collection:
part of the charge of some events is lost near the entrance window or the sensor edges, so
the event is recorded below the full energy.
\blk{Numerics} The tail term overflows when evaluated directly: the exponential diverges
where $\erfc$ underflows. With $z = (x/\sigma+\sigma/\beta)/\sqrt2$ and the scaled
complementary error function $\erfcx(z) = \e^{z^2}\erfc(z)$, the identity
$x/\beta + \sigma^2/2\beta^2 - z^2 = -x^2/2\sigma^2$ gives
\begin{equation}
  \exp\!\left(\frac{x}{\beta}+\frac{\sigma^2}{2\beta^2}\right)\erfc(z) =
  \begin{cases}
    \e^{-x^2/2\sigma^2}\,\erfcx(z), & z>0,\\
    \exp\!\left(\frac{x}{\beta}+\frac{\sigma^2}{2\beta^2}\right)\erfc(z), & z\le0,
  \end{cases}
\end{equation}
both of which are finite (the exponent in the second case is capped at 50).
\begin{hyp}
One tail and one step shared by all lines, in units of $\sigma$; real detectors show a tail
fraction that decreases with energy and may need a second, long tail. The step is a fraction
of the peak \emph{height}, so its area per unit energy scales with $1/\sigma$.
\end{hyp}
\pycode{xrf_fit}{peak_shape}

\subsection{Scatter peaks under monochromatic excitation}\label{sec:scatter}

Under monochromatic excitation the spectrum also contains the photons scattered by the
sample: elastically (Rayleigh) at the excitation energy, and inelastically (Compton) at
lower energy. With a white beam the scattered radiation is a continuum and belongs to the
background.

\paragraph{Compton energy.} A photon of energy $E_0$ scattered by a free electron at rest
through the angle $\theta$ has the energy \citep{Compton1923}
\begin{equation}
  E_C = \frac{E_0}{1 + \dfrac{E_0}{m_ec^2}(1-\cos\theta)} .
  \label{eq:compton}
\end{equation}

\paragraph{Impulse approximation.} Electrons are bound and moving, which broadens the Compton
line. In the impulse approximation \citep{Eisenberger1970,Ribberfors1975} the energy
distribution of photons scattered at angle $\theta$ is
\begin{equation}
  \frac{\dd^2\sigma}{\dd\Omega\,\dd E} \propto J(p_z)\,\frac{\dd p_z}{\dd E},
  \qquad
  p_z = \frac{E_0E(1-\cos\theta)/m_ec^2 - (E_0-E)}{\sqrt{E_0^2+E^2-2E_0E\cos\theta}},
  \label{eq:pz}
\end{equation}
where $p_z$ is the projection of the electron momentum on the scattering vector, in units of
$m_ec$ in Eq.~\eqref{eq:pz} (multiplied by $1/\alpha = 137.036$ to obtain atomic units), and
$J(p_z)$ is the Compton profile of the scatterer: the sum of the tabulated Hartree--Fock
profiles of its atoms \citep{Biggs1975}, available in xraylib as \fn{ComptonProfile}. Outer
electrons give a sharp core, inner shells broad wings. The peak is centred at $p_z=0$, i.e.
at $E_C$.

\paragraph{Range of scattering angles.} A detector close to the sample (2\,mm on the
training spectrum of Section~\ref{sec:notebooks}) collects photons scattered over several
tens of degrees. Since Eq.~\eqref{eq:compton} is non-linear in the angle, the Compton peak
becomes broad and asymmetric; the model averages it over an angle range that is fitted.

\api[\texttt{\_angle\_grid}]{\texttt{\_angle\_grid(angle\_deg, spread\_deg, n=15)}}
Scattering angles sampled uniformly over $\theta\pm\Delta\theta$ (15 points), clipped to
1--179\textdegree; a single angle when the spread is below 0.05\textdegree.
\pycode{xrf_fit}{_angle_grid}

\api[\texttt{ScatterSource}]{\texttt{ScatterSource(energies\_ev, weights=(1.0,), angle\_deg=90.0, scatterer=None)}}
The monochromatic excitation, for the elastic and Compton peaks.
\begin{args}
\item[energies\_ev, weights] excitation lines and their relative intensities (normalised
  internally), e.g. from \fn{tube_lines("Mo")}.
\item[angle\_deg] scattering angle between the incident beam and the direction to the
  detector: the starting value of the fitted Compton angle.
\item[scatterer] chemical formula of the scattering medium (\fn{"H2O"}): the Compton peak then
  has the impulse-approximation shape; without it, the Compton peak is a Gaussian (averaged
  over the angle range).
\end{args}
\blk{Methods} \fn{compton_energy_ev(e_ev, angle_deg)} is Eq.~\eqref{eq:compton};
\fn{profile_table()} tabulates $J(p_z)$ of one formula unit for $0\le p_z\le60$ atomic units
(3001 points), summing the atoms of the formula with their multiplicities; \fn{compton_pz}
is Eq.~\eqref{eq:pz} in atomic units.
\begin{hyp}
\begin{itemize}[nosep,leftmargin=1em]
\item The prefactor of the impulse-approximation cross-section (Klein--Nishina factor and
  $E/E_0$ terms, \citealp{Klein1929,Ribberfors1975}) is taken as constant across the peak, and
  the angular distribution as uniform over $\theta\pm\Delta\theta$ (no Klein--Nishina, solid
  angle or polarisation weighting). The free angle and spread absorb these approximations.
\item The elastic peak has a free amplitude; its intensity is not computed from atomic form
  factors.
\item The scattering medium is the given formula (the matrix), not the full sample; the
  profile of trace elements is negligible.
\end{itemize}
\end{hyp}
\pycode{xrf_fit}{ScatterSource}
\pycode{xrf_fit}{ScatterSource.__post_init__}
\pycode{xrf_fit}{ScatterSource.compton_energy_ev}
\pycode{xrf_fit}{ScatterSource.profile_table}
\pycode{xrf_fit}{ScatterSource.compton_pz}

\subsection{\texttt{XRFModel}}\label{sec:xrfmodel}

\api[\texttt{XRFModel}]{\texttt{XRFModel(groups, calibration, detector, channels, background\_order=1, fit\_resolution=True, fit\_gamma=True, escape=False, pileup\_template=None, fit\_pileup=True, pileup\_tau\_us=0.2, fit\_calibration=False, fit\_air\_path=False, fit\_tails=False, scatter=None, background\_kind="poly")}}
Eq.~\eqref{eq:model} as a callable with the signature that \fn{curve_fit} expects,
\fn{model(channels, *params)}. Table~\ref{tab:params} lists the parameters, in their order in
the parameter vector; \fn{param_names} holds the names.

\begin{table}[h]
\centering\small
\caption{Parameters of \fn{XRFModel}, in order. Optional blocks appear only when the
corresponding option is set.}
\label{tab:params}
\begin{tabularx}{\textwidth}{@{}l l X@{}}
\toprule
Name & Present if & Meaning \\
\midrule
\fn{bkg_0} \dots\ \fn{bkg_n} & \fn{background_order} not \fn{None} & background coefficients
in $u\in[-1,1]$ \\
\fn{A_<group>} & always & reference-line area of each group (counts summed over channels) \\
\fn{noise_ev}, \fn{fano} & \fn{fit_resolution} & resolution law, Eq.~\eqref{eq:resolution} \\
\fn{gamma_ev} & \fn{fit_gamma} & Lorentzian half width \\
\fn{cal_gain}, \fn{cal_offset_ev} & \fn{fit_calibration} & calibration (only the offset with
\fn{"offset"}) \\
\fn{air_path_cm} & \fn{fit_air_path} & sample--detector air path \\
\fn{tail_fraction}, \fn{tail_slope}, \fn{step_fraction} & \fn{fit_tails} & detector line
shape, Eq.~\eqref{eq:hypermet} \\
\fn{A_elastic}, \fn{elastic_sigma_ev} & \fn{scatter} & elastic area and extra Gaussian
width (monochromator bandwidth) \\
\fn{A_compton}, \fn{compton_angle_deg}, \fn{compton_angle_spread_deg} & \fn{scatter} &
Compton area, mean angle and half range \\
\fn{compton_sigma_ev}, \fn{compton_tail_fraction}, \fn{compton_tail_slope} & \fn{scatter} &
extra Gaussian width and low-energy tail of the Compton peak \\
\fn{pileup_tau_us} & template and \fn{fit_pileup} & pulse-pair resolving time (\textmu s) \\
\bottomrule
\end{tabularx}
\end{table}

\begin{args}
\item[groups] from \fn{element_line_groups} (may be empty with a scatter source).
\item[calibration] fixed calibration, or the starting point when it is fitted.
\item[detector] resolution, efficiency and line shape; supplies the starting (or fixed) values
  of the shared parameters.
\item[channels] channel grid of the fit range.
\item[background\_order] polynomial order (\fn{None}: no background term).
\item[background\_kind] \fn{"poly"}: the background is the polynomial; \fn{"exp"}: its
  exponential, positive by construction and closer to a scatter continuum over a wide range
  (non-linear in its coefficients).
\item[fit\_resolution, fit\_gamma] free $N$, $F$ and $\Gamma$.
\item[escape] add the escape peaks of every line, tied to its amplitude (no extra parameter):
  the group amplitude then counts parent and escape, i.e. every photon absorbed in the sensor,
  which is what \fn{Detector.efficiency} describes.
\item[pileup\_template, fit\_pileup, pileup\_tau\_us] pile-up spectrum per \textmu s (full MCA
  range), with $\tau$ free or fixed. The sum peaks of the main lines usually lie above the
  analysis range: fit $\tau$ once over an extended range and fix it in the analysis.
\item[fit\_calibration] \fn{True} (gain and offset), \fn{"offset"} (energy shift only) or
  \fn{False}.
\item[fit\_air\_path] free air path $d$. The relative intensities were computed at $d_0 =$
  \fn{detector.air_path_cm}; at $d$ each line is rescaled by
  $\tau_l = \exp[-(\mu_l-\mu_\mathrm{ref})(d-d_0)]$, with $\mu$ the linear attenuation of air
  at the line and at the group's reference line. The information comes from intensity ratios
  within groups (K$\alpha$/K$\beta$, L$\alpha$/L$\beta$/L$\gamma$), independent of the sample
  composition; the amplitudes absorb the rest.
\item[fit\_tails] free tail and step, shared by all lines.
\item[scatter] a \fn{ScatterSource}: adds the elastic and Compton peaks.
\end{args}
\pycode{xrf_fit}{XRFModel}
\pycode{xrf_fit}{XRFModel.__init__}

The constructor precomputes the escape lines of every line, the air attenuation differences
$\mu_l-\mu_\mathrm{ref}$, the Compton profile table, the energies of the channel grid and the
scaled coordinate $u = 2(c-c_0)/(c_1-c_0) - 1 \in [-1,1]$ in which the background polynomial
is evaluated (powers of $u$ stay of order 1, which keeps the fit well conditioned).

\api{\texttt{\_unpack}, \texttt{calibration\_of}, \texttt{air\_path\_of}, \texttt{pileup\_tau\_us}, \texttt{pileup\_profile}, \texttt{tails\_of}}
Accessors that read a parameter vector: background coefficients, amplitudes and resolution
parameters (fitted values, or the detector's when frozen); the calibration, air path,
resolving time and line-shape parameters, fitted or fixed. \fn{pileup_profile} is
$\tau_\mathrm{pp}P(c)$ on the fit range.
\pycode{xrf_fit}{XRFModel._unpack}
\pycode{xrf_fit}{XRFModel.calibration_of}
\pycode{xrf_fit}{XRFModel.air_path_of}
\pycode{xrf_fit}{XRFModel.pileup_tau_us}
\pycode{xrf_fit}{XRFModel.pileup_profile}
\pycode{xrf_fit}{XRFModel.tails_of}

\api[\texttt{background}, \texttt{\_shape}]{\texttt{background(params)}, \texttt{\_shape(centre\_ch, sigma\_ch, gamma\_ch, tails)}}
$B = \sum_k b_k u^k$, or $\exp(\sum_k b_k u^k)$ with the exponent clipped to $\pm50$.
\fn{_shape} evaluates \fn{peak_shape} on the channel grid.
\pycode{xrf_fit}{XRFModel.background}
\pycode{xrf_fit}{XRFModel._shape}

\api[\texttt{group\_profile}]{\texttt{group\_profile(index, params, sigma\_scale=1.0)}}
The counts of one group, Eq.~\eqref{eq:model} without background: for each line, the parent
peak with weight $A_g r_l\tau_l(1-\sum_e f_{le})$ and each escape peak with weight
$A_g r_l\tau_l f_{le}$. Profiles are evaluated in \emph{channel} space, with widths
$\sigma/G$ in channels: since the profile has unit area over its argument, the amplitude is
directly the number of counts summed over channels. An escape peak has the resolution of its
own (deposited) energy. \fn{sigma_scale} multiplies every Gaussian width; it is used for the
width-correction term of \fn{fit_map}.
\pycode{xrf_fit}{XRFModel.group_profile}

\api[\texttt{scatter\_profile}, \texttt{\_compton\_kernel}, \texttt{\_compton\_ia}]{\texttt{scatter\_profile(params, which="both")}, \texttt{\_compton\_kernel(x, sigma, v, step)}, \texttt{\_compton\_ia(cal, e0, v, noise, fano, tails)}}
The scatter peaks, for each excitation line with its weight:
\begin{itemize}[nosep]
\item \textbf{elastic}: the detector response at $E_0$ with the width
  $\sqrt{\sigma^2(E_0)+\sigma_\mathrm{el}^2}$ (the extra width stands for the bandwidth of the
  monochromator), with its escape peaks;
\item \textbf{Compton}, impulse approximation (with a scatterer): \fn{_compton_ia} evaluates
  $J(|p_z|)\,|\dd p_z/\dd E|$ on a channel grid extended 3\,keV below the fit range (so that
  the wings are not cut), summed over the angle grid, set to zero above $E_0$ and normalised
  to unit area. It is convolved by FFT with the Compton kernel---a Gaussian of width
  $\sqrt{\sigma^2(E_C)+\sigma_C^2}$ with its own low-energy tail
  (\fn{compton_tail_fraction}, \fn{compton_tail_slope})---to which the detector step is added
  as a shelf proportional to the number of counts above each channel. The escape peaks are
  shifted copies of this base profile;
\item \textbf{Compton}, Gaussian (without scatterer): the kernel at $E_C(\theta)$ for each
  angle of the grid, averaged, with escape peaks.
\end{itemize}
Returns zeros without a scatter source.
\pycode{xrf_fit}{XRFModel.scatter_profile}
\pycode{xrf_fit}{XRFModel._compton_kernel}
\pycode{xrf_fit}{XRFModel._compton_ia}

\api[\texttt{\_\_call\_\_}]{\texttt{\_\_call\_\_(channels, *params)}}
The total model: background, pile-up, scatter, and every group.
\pycode{xrf_fit}{XRFModel.__call__}

\api[\texttt{initial\_guess}]{\texttt{initial\_guess(y)}}
Starting values and bounds from the data. The background starts at the 5th percentile of the
spectrum (its logarithm for \fn{"exp"}), other coefficients at 0. Each amplitude is seeded
from the observed height above that baseline at the strongest line of the group inside the
range, divided by its relative intensity and converted to an area with the modelled width,
$A = h\,\sigma\sqrt{2\pi}$. Bounds: amplitudes $\ge0$; noise 1--200\,eV; Fano 0.01--1;
$\Gamma$ 0--100\,eV; gain $\pm20\,\%$; offset $\pm50$ channels; air path 0--200\,cm; tail
fraction 0--0.6, slope 0.2--20\,$\sigma$, step 0--0.1; Compton angle 10--175\textdegree,
spread 0--80\textdegree, extra widths up to 0.5\,keV (elastic) and 2\,keV (Compton); resolving
time 0--50\,\textmu s.
\begin{hyp}
The non-linear parameters (calibration, Compton angle) need a good start: a calibration off
by tens of channels can lead to a wrong local minimum. Start from a calibration fitted on
known peaks (\fn{xrf_calibration.calibration_from_peaks} or \fn{calibrate}).
\end{hyp}
\pycode{xrf_fit}{XRFModel.initial_guess}

\subsection{Fitting a spectrum}\label{sec:fitting}

\api{\texttt{FitResult}}
The outcome of a fit: the \fn{model}, the parameters \fn{popt}, their errors \fn{perr} and
covariance \fn{pcov}, the fitted \fn{channels}, data \fn{y}, errors \fn{yerr} and model
\fn{yfit}, the reduced $\chi^2$, and the group \fn{areas} and \fn{area_errors} (reference-line
areas). Properties and methods: \fn{calibration} (fitted or fixed), \fn{residual},
\fn{residual_normalised} (in units of the data error), \fn{params()} (\texttt{\{name: (value,
error)\}}), \fn{summary()} (text report with the total area of every group).
\pycode{xrf_fit}{FitResult}
\pycode{xrf_fit}{FitResult.calibration}
\pycode{xrf_fit}{FitResult.residual}
\pycode{xrf_fit}{FitResult.residual_normalised}
\pycode{xrf_fit}{FitResult.params}
\pycode{xrf_fit}{FitResult.summary}

\api[\texttt{fit\_average\_spectrum}]{\texttt{fit\_average\_spectrum(y, groups, calibration, detector, roi, yerr=None, background\_order=1, fit\_resolution=True, fit\_gamma=True, maxfev=20000, escape=False, pileup\_template=None, fit\_pileup=True, pileup\_tau\_us=0.2, fit\_calibration=False, fit\_air\_path=False, fit\_tails=False, scatter=None, background\_kind="poly")}}
Fits Eq.~\eqref{eq:model} to one spectrum (a map average, a single point, a reference
spectrum).
\begin{args}
\item[y] spectrum over the full MCA range (sliced to \fn{roi} internally).
\item[groups, calibration, detector] as for \fn{XRFModel}.
\item[roi] half-open channel range of the fit.
\item[yerr] one-sigma errors per channel (e.g. \fn{counting_error}); zero errors are raised to
  the smallest positive one, which would otherwise make the fit singular.
\item[maxfev] limit on function evaluations.
\item[others] passed to \fn{XRFModel}.
\end{args}
\blk{Method} Weighted non-linear least squares,
\begin{equation}
  \chi^2(\mathbf p) = \sum_c \left(\frac{y_c - y(c;\mathbf p)}{\sigma_c}\right)^2,
\end{equation}
minimised by \fn{scipy.optimize.curve_fit} with bounds, i.e.\ the trust-region reflective
algorithm \citep{Branch1999,Virtanen2020}. With \fn{absolute_sigma=True} the covariance is
$(\mathbf J^\mathsf T\mathbf W\mathbf J)^{-1}$ at the minimum, not rescaled by $\chi^2_\nu$, so
the reported errors are the statistical ones \citep{Bevington2003}. The reduced
$\chi^2_\nu = \chi^2/(n-p)$ measures the quality of the model: values well above 1 mean that
the residuals are dominated by model errors, and the statistical errors then underestimate
the true uncertainty.
\begin{hyp}
Gaussian errors. Minimising $\chi^2$ with data-derived weights (Neyman's $\chi^2$) biases
the amplitudes low at low counts \citep{Baker1984,Mighell1999}; map averages and long single
measurements have many counts per channel, where the bias is negligible.
\end{hyp}
\blk{Returns} a \fn{FitResult}.
\pycode{xrf_fit}{fit_average_spectrum}

\subsection{Maps: the per-pixel linear solve}

Once the shape parameters (resolution, calibration, tails, pile-up time) are fixed by the fit
of the average spectrum, Eq.~\eqref{eq:model} is \emph{linear} in the remaining parameters:
the background coefficients and the group amplitudes. Every pixel is then solved exactly and
at once, with its own covariance, instead of running a non-linear fit per pixel.

\api{\texttt{MapResult}}
Per-pixel results: \fn{group_labels}, \fn{areas} and \fn{area_errors} $(n_\mathrm{pixels},
n_\mathrm{groups})$, \fn{background} coefficients, \fn{chi2_red}, \fn{total_rel}, map
\fn{shape}, and optionally \fn{width_change} and \fn{segments}.
\blk{Methods}
\fn{area_map(label, total=False)} and \fn{error_map} reshape a group's areas (or errors) to
the map, optionally scaled to the whole series; \fn{snr} is their ratio;
\fn{width_change_map} reshapes the width term; \fn{ratio_map(numerator, denominator)} returns
the ratio of two groups' areas with the propagated error,
$\sigma_r = |r|\sqrt{(\sigma_a/a)^2+(\sigma_b/b)^2}$; \fn{quality_report} lists, per group, the
median area, the median signal-to-noise and the fraction of negative pixels, and flags groups
as quantitative (median S/N $\ge$ 20 by default), detected only (S/N $\ge$ 3: the map is noise
and the average area a lower bound), or not detected.
\pycode{xrf_fit}{MapResult}
\pycode{xrf_fit}{MapResult.width_change_map}
\pycode{xrf_fit}{MapResult.area_map}
\pycode{xrf_fit}{MapResult.error_map}
\pycode{xrf_fit}{MapResult.snr}
\pycode{xrf_fit}{MapResult.quality_report}
\pycode{xrf_fit}{MapResult.ratio_map}

\api[\texttt{number\_ratio}]{\texttt{number\_ratio(area\_a, area\_b, group\_a, group\_b, err\_a=None, err\_b=None)}}
Atom-number ratio from two reference-line areas,
\begin{equation}
  \frac{n_a}{n_b} = \frac{A_a}{A_b}\,\frac{S_bW_b}{S_aW_a},
\end{equation}
with $S$ the groups' \fn{ref_strength} (mass cross-section $\times$ excitation $\times$
efficiency) and $W$ the atomic weights, which turn mass cross-sections into per-atom ones.
Errors are propagated as uncorrelated. Returns \fn{(ratio, ratio_err, k)} with $k$ the factor
multiplying $A_a/A_b$.
\begin{hyp}
Both groups come from the same \fn{element_line_groups} configuration, and matrix effects
are negligible or equal for the two lines (thin sample). Otherwise use
\fn{xrf_matrix.quantify_layer}.
\end{hyp}
\pycode{xrf_fit}{number_ratio}

\api[\texttt{design\_matrix}]{\texttt{design\_matrix(model, shape\_params)}}
The columns of the linear model: $u^k$ for the background terms, then the profile of each
group with unit amplitude (\fn{group_profile} with all other amplitudes set to 0). Returns an
$(n_\mathrm{channels}, n_\mathrm{bkg}+n_\mathrm{groups})$ array.
\pycode{xrf_fit}{design_matrix}

\api[\texttt{fit\_map}]{\texttt{fit\_map(raw\_spectra, scale, model, shape\_params, map\_shape, poisson\_weights=True, clip\_negative=False, pileup=None, segments=None, width\_correction=False)}}
Solves the fixed-shape model for every pixel.
\begin{args}
\item[raw\_spectra] raw counts $(n_\mathrm{pixels}, n_\mathrm{channels})$: the weights are
  defined on counts, and the normalisation is applied internally.
\item[scale] the normalisation divisors of \fn{normalise_spectra}.
\item[model] an \fn{XRFModel} for the fit range (no scatter, polynomial background).
\item[shape\_params] parameter vector of the average fit, or \texttt{\{segment: vector\}}.
\item[map\_shape] $(n_\mathrm{slow}, n_\mathrm{fast})$.
\item[poisson\_weights] weight each channel by $1/\mathrm{Var}(y) = s^2/\max(x,1)$. Without
  them the strongest channels dominate every pixel and the errors are meaningless.
\item[clip\_negative] clip negative amplitudes to zero (off: a negative amplitude is a
  diagnostic of a mis-specified model).
\item[pileup] per-pixel pile-up spectra; with a pile-up term in the model,
  $\tau P_i/s_i$ is subtracted from each pixel before the solve (a known, pixel-specific
  component).
\item[segments] per-pixel segment labels, with \fn{shape_params} per segment: each segment
  is solved with its own widths, the exact treatment of an abrupt resolution change.
\item[width\_correction] add one first-order width term per pixel (below).
\end{args}
\blk{Method} For each pixel $p$ with design matrix $\mathbf A$ and weights $\mathbf W_p$, the
weighted least-squares solution and its covariance are
\begin{equation}
  \hat{\mathbf c}_p = (\mathbf A^\mathsf T\mathbf W_p\mathbf A)^{-1}\mathbf A^\mathsf T
  \mathbf W_p\mathbf y_p, \qquad
  \mathrm{Cov}(\hat{\mathbf c}_p) = (\mathbf A^\mathsf T\mathbf W_p\mathbf A)^{-1},
\end{equation}
assembled for all pixels at once with \fn{einsum} (each system is only $k\times k$) and solved
in one batched call; a relative ridge of $10^{-12}$ times the trace keeps empty pixels from
making a system singular (they are set to NaN afterwards). The reduced $\chi^2$ of each pixel
uses the same weights.
\blk{Width correction} If the detector resolution changes slightly during the scan, the fixed
widths bias overlapping peaks. The extra column is the derivative of the modelled peaks with
respect to a common relative scaling $\varepsilon$ of every Gaussian width, at the shape
amplitudes $\mathbf a$, computed by central differences ($\pm2\,\%$). A pixel modelled as
$\lambda\sum_j a_j g_j[(1+\varepsilon)\sigma] \approx \lambda\sum_j a_j g_j + \lambda\varepsilon
\,\partial_\varepsilon(\cdot)$ has a width coefficient $\delta = \lambda\varepsilon$, with
$\lambda$ the projection of the pixel's amplitudes on $\mathbf a$; hence $\varepsilon =
\delta/\lambda$ is returned as \fn{width_change}. First order: valid for changes of a few per
cent, and if the relative amplitudes stay close to those of the average.
\pycode{xrf_fit}{fit_map}

\subsection{Resolution stability during a scan}

\api[\texttt{width\_window}]{\texttt{width\_window(groups, areas, detector, min\_frac=0.03, doublet\_ev=60.0)}}
Chooses a clean energy window around a strong, isolated line to monitor the resolution. For
each group's reference line, the half width is $\min(2.5\,\mathrm{FWHM}, d/2)$ with $d$ the
distance to the nearest line of expected area above \fn{min_frac} of the candidate's (lines
closer than \fn{doublet_ev}, such as K$\alpha_{1,2}$, belong to the peak); the score is the
area times the Gaussian fraction inside the window, $\mathrm{erf}(h/\sqrt2\sigma)$, and
windows narrower than 1.2\,FWHM are rejected. Returns \fn{(emin_ev, emax_ev, label)}.
\pycode{xrf_fit}{width_window}

\api[\texttt{peak\_width\_profile}]{\texttt{peak\_width\_profile(raw\_spectra, calibration, window\_ev, bin\_size)}}
The centroid and r.m.s. width of a peak along the acquisition order. Consecutive pixels are
summed in bins (one map row is a natural bin), a linear background between the window edges
is removed, and the first and second moments of the net peak give the centroid $\bar E$ and
the width $w$, with the statistical error $w/\sqrt{2n}$ of a Gaussian width estimated from
$n$ counts. A resolution change shows as a step in $w$, a gain or offset change as a step in
$\bar E$. Independent of the normalisation. Returns \texttt{\{width\_ev, centroid\_ev, counts,
width\_err\_ev\}}.
\pycode{xrf_fit}{peak_width_profile}

\api[\texttt{detect\_steps}]{\texttt{detect\_steps(values, window=5, threshold=6.0, min\_step=0.0)}}
Finds abrupt jumps in a series. At each index the difference between the means of the
\fn{window} points after and before is divided by its noise, $\sigma\sqrt{2/w}$, where
$\sigma$ is the point-to-point noise estimated robustly from the median absolute deviation of
successive differences, $\sigma = 1.4826\,\mathrm{MAD}(\Delta v)/\sqrt2$ \citep{Hampel1974}.
Local maxima of this statistic above \fn{threshold}, with a step larger than \fn{min_step},
are returned (a new level starts at each index). Slow drifts are better handled by the width
term of \fn{fit_map}.
\pycode{xrf_fit}{detect_steps}

\api[\texttt{segment\_labels}, \texttt{fit\_segments}]{\texttt{segment\_labels(steps, n\_bins, bin\_size, n\_pixels=None)}, \texttt{fit\_segments(corrected, raw\_spectra, scale, segments, groups, calibration, detector, roi, pileup=None, **fit\_kw)}}
\fn{segment_labels} turns step positions (in bins) into a label per pixel (0, 1, \dots), the
last incomplete bin taking the last label. \fn{fit_segments} fits the average spectrum of
each segment, with its own counting error and pile-up template, and returns
\texttt{\{label: FitResult\}}; \texttt{\{label: r.popt\}} then goes to \fn{fit_map} with the same
segments.
\pycode{xrf_fit}{segment_labels}
\pycode{xrf_fit}{fit_segments}

\subsection{Diagnostics and plots}

\api[\texttt{identify\_peaks}]{\texttt{identify\_peaks(y, calibration, roi, prominence\_frac=0.01, tol\_ev=70.0, min\_rate=0.08, z\_range=(12, 83), n\_candidates=4)}}
A scouting tool for choosing the element list: peaks are found with
\fn{scipy.signal.find_peaks} (prominence above \fn{prominence_frac} of the maximum, at least 8
channels apart), and for each peak the tabulated IUPAC lines within \fn{tol_ev}, of elements
in \fn{z_range} and with a radiative rate above \fn{min_rate}, are listed, strongest rate
first. An energy coincidence is only a hint: an element is confirmed when its whole line
series fits. Returns \texttt{[\{channel, energy\_ev, height, candidates\}]}.
\pycode{xrf_fit}{identify_peaks}

\api[\texttt{pixel\_parameters}, \texttt{select\_pixels}]{\texttt{pixel\_parameters(model, map\_result, shape\_params, pixel)}, \texttt{select\_pixels(map\_result, n=10, mode="random", seed=0)}}
\fn{pixel_parameters} rebuilds a full parameter vector for one pixel (its background and
amplitudes from the map solve, the shape from the average fit or its segment).
\fn{select_pixels} picks pixel indices to inspect: \fn{"random"} among finite fits,
\fn{"worst"} and \fn{"best"} by reduced $\chi^2$, or \fn{"spread"} along the $\chi^2$
distribution in equal quantiles.
\pycode{xrf_fit}{pixel_parameters}
\pycode{xrf_fit}{select_pixels}

\api[\texttt{plot\_pixel\_fits}]{\texttt{plot\_pixel\_fits(raw\_spectra, scale, model, map\_result, shape\_params, pixels=None, n=10, mode="random", seed=0, ncols=5, energy\_axis=True, logy=True, show\_components=True, figsize\_per\_panel=(3.2, 2.6), pileup=None)}}
A grid of individual pixel spectra with their fitted model, components shaded and the
pixel's $\chi^2_\nu$ in the title: the direct check that the fixed line ratios hold pixel by
pixel and that the background is plausible at the local count level. With pile-up, the
per-pixel pile-up spectra replace the map average. Returns \fn{(fig, axes, pixels)}.
\pycode{xrf_fit}{plot_pixel_fits}

\api[\texttt{plot\_fit}]{\texttt{plot\_fit(result, energy\_axis=True, logy=False, figsize=(11.0, 7.0), data\_label="map-averaged data", ylabel=...)}}
Two panels: the spectrum with its error band, the background, every group (shaded, with its
reference line marked), the scatter peaks if any, and the total model; below, the residual in
units of the data error. Returns \fn{(fig, (ax_spectrum, ax_residual))}.
\pycode{xrf_fit}{plot_fit}
